\begin{abstract}
A handful of observed responses can sharpen forecasts of an AI system's future performance. We present \modelname{}, a compact evaluator that combines two complementary views of adaptation. A Rasch branch updates the system's overall ability; a two-dimensional CAIMIRA-style branch determines how the update transfers across items. Both start from item text and subject metadata, and their probabilities receive equal weight. The CAIMIRA update decomposes into a shared shift and a content-dependent contrast, making the route from revealed evidence to new forecasts explicit. Our selected PAIEC candidate, \code{18-s44-rasch-caimira.zip}, achieves formative Brier area under the learning curve of \textbf{0.160822}. Across nine subject--benchmark pairs, mean Brier error falls from 0.212061 with zero labels to 0.126671 with 31 labels, a \textbf{40.27\% reduction}. Every pair improves at the final budget, and seven labels capture 76.20\% of the observed endpoint gain. The archive contains 3,076 fitted coefficients and runs with local CPU computation. We explain its adaptation geometry, analyze the complete learning curve, and develop hierarchical evidence sharing and calibrated gating as complementary extensions.
\end{abstract}

\section{Introduction}
Predictive evaluation turns a small evaluation budget into forecasts for a much larger set of items. Rather than treating a subject's observed accuracy as the final answer, it asks how each revealed outcome changes the probability of success on further questions. This creates a useful division of labor: item text and system metadata describe the problem before any labels arrive, while observed responses adapt that description to the current subject and benchmark~\citep{paiec2026}.

The central conceptual step in our design is to treat adaptation at two scales. A response can revise a subject's \emph{overall level}, and it can revise the \emph{profile of ability} that matters for a particular item. The Rasch branch estimates a shared scalar correction. The CAIMIRA-style branch estimates both a shared shift and an item-dependent contrast through two ability coordinates and item-specific relevance weights. Together, these views explain how evidence from one item changes the forecast for another.

This perspective joins the stability of classical item-response modeling~\citep{rasch1960} with CAIMIRA's content-dependent difficulty and relevance~\citep{gor2024}. A fixed lexical representation keeps the fitted model compact~\citep{weinberger2009}. At evaluation time, only the subject's ability is updated. The item model stays fixed, so each new label has a precise route into later predictions. Figure~\ref{fig:overview} shows this computation.

\paragraph{Contributions.}
First, we give a concrete formulation of global and item-sensitive adaptation in one deployable evaluator. Second, we show that the selected candidate converts sparse feedback into substantial probability-accuracy gains: every evaluated pair improves from zero to 31 labels, with most of the aggregate gain available after seven. Third, we connect the model's transparent update rule to a practical extension based on shared item structure, subject regression, and calibrated gating.

